% Shared macros for PYNE arXiv paper series
% Include with: % Shared macros for PYNE arXiv paper series
% Include with: % Shared macros for PYNE arXiv paper series
% Include with: % Shared macros for PYNE arXiv paper series
% Include with: \input{../common/macros.tex}

\usepackage{amsmath,amssymb,amsfonts}
\usepackage{amsthm}
\usepackage{booktabs}
\usepackage{graphicx}
\usepackage{xcolor}
\usepackage{hyperref}
\usepackage{url}
\usepackage{algorithm}
\usepackage{algpseudocode}
\usepackage{listings}
\usepackage{tikz}
\usepackage{pgfplots}
\pgfplotsset{compat=1.17}
\usetikzlibrary{arrows.meta,positioning,fit,calc,shapes.geometric,shapes.multipart,decorations.pathreplacing,backgrounds,chains,matrix,patterns}
\usepackage{subcaption}
\usepackage{multirow}
\usepackage{tabularx}
\usepackage{enumitem}
\usepackage{microtype}
\usepackage{balance}
\usepackage{csquotes}
\usepackage{natbib}

% Colors
\definecolor{pyneorange}{HTML}{F97316}
\definecolor{pyneblue}{HTML}{2563EB}
\definecolor{pynegray}{HTML}{6B7280}
\definecolor{pynegreen}{HTML}{16A34A}
\definecolor{codebg}{HTML}{F8FAFC}
\definecolor{codekw}{HTML}{7C3AED}
\definecolor{codestr}{HTML}{B45309}
\definecolor{codecomment}{HTML}{64748B}

\hypersetup{
  colorlinks=true,
  linkcolor=pyneblue,
  citecolor=pynegreen,
  urlcolor=pyneorange,
  pdftitle={},
  pdfauthor={HOOX / PYNE Project}
}

% Code listings
\lstdefinestyle{pine}{
  basicstyle=\ttfamily\footnotesize,
  backgroundcolor=\color{codebg},
  keywordstyle=\color{codekw}\bfseries,
  stringstyle=\color{codestr},
  commentstyle=\color{codecomment}\itshape,
  breaklines=true,
  frame=single,
  rulecolor=\color{black!20},
  numbers=left,
  numberstyle=\tiny\color{pynegray},
  numbersep=6pt,
  xleftmargin=1.5em,
  framexleftmargin=1.2em,
  showstringspaces=false,
  tabsize=2,
  morekeywords={indicator,strategy,library,plot,plotshape,fill,hline,var,varip,if,else,for,while,switch,true,false,na,and,or,not,import,export,type,method,enum},
  morecomment=[l]{//},
  morestring=[b]",
}
\lstdefinestyle{python}{
  language=Python,
  basicstyle=\ttfamily\footnotesize,
  backgroundcolor=\color{codebg},
  keywordstyle=\color{codekw}\bfseries,
  stringstyle=\color{codestr},
  commentstyle=\color{codecomment}\itshape,
  breaklines=true,
  frame=single,
  rulecolor=\color{black!20},
  numbers=left,
  numberstyle=\tiny\color{pynegray},
  numbersep=6pt,
  xleftmargin=1.5em,
  framexleftmargin=1.2em,
  showstringspaces=false,
  tabsize=4,
}
\lstset{style=python}

% Theorem environments
\theoremstyle{plain}
\newtheorem{theorem}{Theorem}[section]
\newtheorem{lemma}[theorem]{Lemma}
\newtheorem{proposition}[theorem]{Proposition}
\newtheorem{corollary}[theorem]{Corollary}
\theoremstyle{definition}
\newtheorem{definition}[theorem]{Definition}
\newtheorem{example}[theorem]{Example}
\newtheorem{remark}[theorem]{Remark}

% Shortcuts
\newcommand{\pyne}{\textsc{Pyne}}
\newcommand{\pine}{Pine Script\texttrademark}
\newcommand{\tv}{TradingView\textregistered}
\newcommand{\asdl}{\textsc{Asdl}}
\newcommand{\antlr}{\textsc{Antlr}4}
\newcommand{\numba}{\textsc{Numba}}
\newcommand{\lsp}{\textsc{Lsp}}
\newcommand{\ohlcv}{\textsc{Ohlcv}}
\newcommand{\udt}{\textsc{Udt}}
\newcommand{\jit}{\textsc{Jit}}
\newcommand{\api}{\textsc{Api}}
\newcommand{\cli}{\textsc{Cli}}
% AST acronym (text mode). Keep amsmath's math asterisk \ast intact.
\DeclareRobustCommand{\Ast}{\textsc{Ast}}
% Papers in this series write \ast{} for the AST acronym in prose.
% Redefine robustly; use \mathast for a binary asterisk in formulas.
\let\mathast\ast
\DeclareRobustCommand{\ast}{\textsc{Ast}}
\newcommand{\ir}{\textsc{Ir}}
\newcommand{\ta}{\texttt{ta}}
\newcommand{\code}[1]{\texttt{#1}}
\newcommand{\file}[1]{\texttt{#1}}
\newcommand{\eg}{e.g.}
\newcommand{\ie}{i.e.}
\newcommand{\etal}{\textit{et~al.}}
\newcommand{\resp}{respectively}

% Math helpers
\newcommand{\R}{\mathbb{R}}
\newcommand{\N}{\mathbb{N}}
\newcommand{\Z}{\mathbb{Z}}
% Text-mode NA token for prose; use $\namath$ inside formulas if needed.
\DeclareRobustCommand{\na}{\textsf{na}}
\newcommand{\namath}{\mathsf{na}}
\newcommand{\series}[1]{\mathbf{#1}}
\newcommand{\baridx}{i}
\newcommand{\bars}{n}
\newcommand{\period}{p}
\newcommand{\abserr}{\varepsilon_{\mathrm{abs}}}
\newcommand{\relerr}{\varepsilon_{\mathrm{rel}}}
\newcommand{\maxerr}{\varepsilon_{\mathrm{max}}}
\newcommand{\meanerr}{\varepsilon_{\mathrm{mean}}}

% Pipeline notation — use inside math: $\pipeline$
\newcommand{\pipeline}{\mathrm{src}\!\to\!\mathrm{lex}\!\to\!\mathrm{parse}\!\to\!\mathrm{AST}\!\to\!\mathrm{eval}}

% Disclaimer footnote helper
\newcommand{\trademarknote}{%
  \pine{} and \tv{} are trademarks of TradingView, Inc.
  \pyne{} is an independent, unofficial implementation and is not affiliated with,
  authorized by, sponsored by, or endorsed by TradingView, Inc.
}

% Figure path helper (papers live one level under .papers/)
\graphicspath{{figures/}{../common/figures/}}


\usepackage{amsmath,amssymb,amsfonts}
\usepackage{amsthm}
\usepackage{booktabs}
\usepackage{graphicx}
\usepackage{xcolor}
\usepackage{hyperref}
\usepackage{url}
\usepackage{algorithm}
\usepackage{algpseudocode}
\usepackage{listings}
\usepackage{tikz}
\usepackage{pgfplots}
\pgfplotsset{compat=1.17}
\usetikzlibrary{arrows.meta,positioning,fit,calc,shapes.geometric,shapes.multipart,decorations.pathreplacing,backgrounds,chains,matrix,patterns}
\usepackage{subcaption}
\usepackage{multirow}
\usepackage{tabularx}
\usepackage{enumitem}
\usepackage{microtype}
\usepackage{balance}
\usepackage{csquotes}
\usepackage{natbib}

% Colors
\definecolor{pyneorange}{HTML}{F97316}
\definecolor{pyneblue}{HTML}{2563EB}
\definecolor{pynegray}{HTML}{6B7280}
\definecolor{pynegreen}{HTML}{16A34A}
\definecolor{codebg}{HTML}{F8FAFC}
\definecolor{codekw}{HTML}{7C3AED}
\definecolor{codestr}{HTML}{B45309}
\definecolor{codecomment}{HTML}{64748B}

\hypersetup{
  colorlinks=true,
  linkcolor=pyneblue,
  citecolor=pynegreen,
  urlcolor=pyneorange,
  pdftitle={},
  pdfauthor={HOOX / PYNE Project}
}

% Code listings
\lstdefinestyle{pine}{
  basicstyle=\ttfamily\footnotesize,
  backgroundcolor=\color{codebg},
  keywordstyle=\color{codekw}\bfseries,
  stringstyle=\color{codestr},
  commentstyle=\color{codecomment}\itshape,
  breaklines=true,
  frame=single,
  rulecolor=\color{black!20},
  numbers=left,
  numberstyle=\tiny\color{pynegray},
  numbersep=6pt,
  xleftmargin=1.5em,
  framexleftmargin=1.2em,
  showstringspaces=false,
  tabsize=2,
  morekeywords={indicator,strategy,library,plot,plotshape,fill,hline,var,varip,if,else,for,while,switch,true,false,na,and,or,not,import,export,type,method,enum},
  morecomment=[l]{//},
  morestring=[b]",
}
\lstdefinestyle{python}{
  language=Python,
  basicstyle=\ttfamily\footnotesize,
  backgroundcolor=\color{codebg},
  keywordstyle=\color{codekw}\bfseries,
  stringstyle=\color{codestr},
  commentstyle=\color{codecomment}\itshape,
  breaklines=true,
  frame=single,
  rulecolor=\color{black!20},
  numbers=left,
  numberstyle=\tiny\color{pynegray},
  numbersep=6pt,
  xleftmargin=1.5em,
  framexleftmargin=1.2em,
  showstringspaces=false,
  tabsize=4,
}
\lstset{style=python}

% Theorem environments
\theoremstyle{plain}
\newtheorem{theorem}{Theorem}[section]
\newtheorem{lemma}[theorem]{Lemma}
\newtheorem{proposition}[theorem]{Proposition}
\newtheorem{corollary}[theorem]{Corollary}
\theoremstyle{definition}
\newtheorem{definition}[theorem]{Definition}
\newtheorem{example}[theorem]{Example}
\newtheorem{remark}[theorem]{Remark}

% Shortcuts
\newcommand{\pyne}{\textsc{Pyne}}
\newcommand{\pine}{Pine Script\texttrademark}
\newcommand{\tv}{TradingView\textregistered}
\newcommand{\asdl}{\textsc{Asdl}}
\newcommand{\antlr}{\textsc{Antlr}4}
\newcommand{\numba}{\textsc{Numba}}
\newcommand{\lsp}{\textsc{Lsp}}
\newcommand{\ohlcv}{\textsc{Ohlcv}}
\newcommand{\udt}{\textsc{Udt}}
\newcommand{\jit}{\textsc{Jit}}
\newcommand{\api}{\textsc{Api}}
\newcommand{\cli}{\textsc{Cli}}
% AST acronym (text mode). Keep amsmath's math asterisk \ast intact.
\DeclareRobustCommand{\Ast}{\textsc{Ast}}
% Papers in this series write \ast{} for the AST acronym in prose.
% Redefine robustly; use \mathast for a binary asterisk in formulas.
\let\mathast\ast
\DeclareRobustCommand{\ast}{\textsc{Ast}}
\newcommand{\ir}{\textsc{Ir}}
\newcommand{\ta}{\texttt{ta}}
\newcommand{\code}[1]{\texttt{#1}}
\newcommand{\file}[1]{\texttt{#1}}
\newcommand{\eg}{e.g.}
\newcommand{\ie}{i.e.}
\newcommand{\etal}{\textit{et~al.}}
\newcommand{\resp}{respectively}

% Math helpers
\newcommand{\R}{\mathbb{R}}
\newcommand{\N}{\mathbb{N}}
\newcommand{\Z}{\mathbb{Z}}
% Text-mode NA token for prose; use $\namath$ inside formulas if needed.
\DeclareRobustCommand{\na}{\textsf{na}}
\newcommand{\namath}{\mathsf{na}}
\newcommand{\series}[1]{\mathbf{#1}}
\newcommand{\baridx}{i}
\newcommand{\bars}{n}
\newcommand{\period}{p}
\newcommand{\abserr}{\varepsilon_{\mathrm{abs}}}
\newcommand{\relerr}{\varepsilon_{\mathrm{rel}}}
\newcommand{\maxerr}{\varepsilon_{\mathrm{max}}}
\newcommand{\meanerr}{\varepsilon_{\mathrm{mean}}}

% Pipeline notation — use inside math: $\pipeline$
\newcommand{\pipeline}{\mathrm{src}\!\to\!\mathrm{lex}\!\to\!\mathrm{parse}\!\to\!\mathrm{AST}\!\to\!\mathrm{eval}}

% Disclaimer footnote helper
\newcommand{\trademarknote}{%
  \pine{} and \tv{} are trademarks of TradingView, Inc.
  \pyne{} is an independent, unofficial implementation and is not affiliated with,
  authorized by, sponsored by, or endorsed by TradingView, Inc.
}

% Figure path helper (papers live one level under .papers/)
\graphicspath{{figures/}{../common/figures/}}


\usepackage{amsmath,amssymb,amsfonts}
\usepackage{amsthm}
\usepackage{booktabs}
\usepackage{graphicx}
\usepackage{xcolor}
\usepackage{hyperref}
\usepackage{url}
\usepackage{algorithm}
\usepackage{algpseudocode}
\usepackage{listings}
\usepackage{tikz}
\usepackage{pgfplots}
\pgfplotsset{compat=1.17}
\usetikzlibrary{arrows.meta,positioning,fit,calc,shapes.geometric,shapes.multipart,decorations.pathreplacing,backgrounds,chains,matrix,patterns}
\usepackage{subcaption}
\usepackage{multirow}
\usepackage{tabularx}
\usepackage{enumitem}
\usepackage{microtype}
\usepackage{balance}
\usepackage{csquotes}
\usepackage{natbib}

% Colors
\definecolor{pyneorange}{HTML}{F97316}
\definecolor{pyneblue}{HTML}{2563EB}
\definecolor{pynegray}{HTML}{6B7280}
\definecolor{pynegreen}{HTML}{16A34A}
\definecolor{codebg}{HTML}{F8FAFC}
\definecolor{codekw}{HTML}{7C3AED}
\definecolor{codestr}{HTML}{B45309}
\definecolor{codecomment}{HTML}{64748B}

\hypersetup{
  colorlinks=true,
  linkcolor=pyneblue,
  citecolor=pynegreen,
  urlcolor=pyneorange,
  pdftitle={},
  pdfauthor={HOOX / PYNE Project}
}

% Code listings
\lstdefinestyle{pine}{
  basicstyle=\ttfamily\footnotesize,
  backgroundcolor=\color{codebg},
  keywordstyle=\color{codekw}\bfseries,
  stringstyle=\color{codestr},
  commentstyle=\color{codecomment}\itshape,
  breaklines=true,
  frame=single,
  rulecolor=\color{black!20},
  numbers=left,
  numberstyle=\tiny\color{pynegray},
  numbersep=6pt,
  xleftmargin=1.5em,
  framexleftmargin=1.2em,
  showstringspaces=false,
  tabsize=2,
  morekeywords={indicator,strategy,library,plot,plotshape,fill,hline,var,varip,if,else,for,while,switch,true,false,na,and,or,not,import,export,type,method,enum},
  morecomment=[l]{//},
  morestring=[b]",
}
\lstdefinestyle{python}{
  language=Python,
  basicstyle=\ttfamily\footnotesize,
  backgroundcolor=\color{codebg},
  keywordstyle=\color{codekw}\bfseries,
  stringstyle=\color{codestr},
  commentstyle=\color{codecomment}\itshape,
  breaklines=true,
  frame=single,
  rulecolor=\color{black!20},
  numbers=left,
  numberstyle=\tiny\color{pynegray},
  numbersep=6pt,
  xleftmargin=1.5em,
  framexleftmargin=1.2em,
  showstringspaces=false,
  tabsize=4,
}
\lstset{style=python}

% Theorem environments
\theoremstyle{plain}
\newtheorem{theorem}{Theorem}[section]
\newtheorem{lemma}[theorem]{Lemma}
\newtheorem{proposition}[theorem]{Proposition}
\newtheorem{corollary}[theorem]{Corollary}
\theoremstyle{definition}
\newtheorem{definition}[theorem]{Definition}
\newtheorem{example}[theorem]{Example}
\newtheorem{remark}[theorem]{Remark}

% Shortcuts
\newcommand{\pyne}{\textsc{Pyne}}
\newcommand{\pine}{Pine Script\texttrademark}
\newcommand{\tv}{TradingView\textregistered}
\newcommand{\asdl}{\textsc{Asdl}}
\newcommand{\antlr}{\textsc{Antlr}4}
\newcommand{\numba}{\textsc{Numba}}
\newcommand{\lsp}{\textsc{Lsp}}
\newcommand{\ohlcv}{\textsc{Ohlcv}}
\newcommand{\udt}{\textsc{Udt}}
\newcommand{\jit}{\textsc{Jit}}
\newcommand{\api}{\textsc{Api}}
\newcommand{\cli}{\textsc{Cli}}
% AST acronym (text mode). Keep amsmath's math asterisk \ast intact.
\DeclareRobustCommand{\Ast}{\textsc{Ast}}
% Papers in this series write \ast{} for the AST acronym in prose.
% Redefine robustly; use \mathast for a binary asterisk in formulas.
\let\mathast\ast
\DeclareRobustCommand{\ast}{\textsc{Ast}}
\newcommand{\ir}{\textsc{Ir}}
\newcommand{\ta}{\texttt{ta}}
\newcommand{\code}[1]{\texttt{#1}}
\newcommand{\file}[1]{\texttt{#1}}
\newcommand{\eg}{e.g.}
\newcommand{\ie}{i.e.}
\newcommand{\etal}{\textit{et~al.}}
\newcommand{\resp}{respectively}

% Math helpers
\newcommand{\R}{\mathbb{R}}
\newcommand{\N}{\mathbb{N}}
\newcommand{\Z}{\mathbb{Z}}
% Text-mode NA token for prose; use $\namath$ inside formulas if needed.
\DeclareRobustCommand{\na}{\textsf{na}}
\newcommand{\namath}{\mathsf{na}}
\newcommand{\series}[1]{\mathbf{#1}}
\newcommand{\baridx}{i}
\newcommand{\bars}{n}
\newcommand{\period}{p}
\newcommand{\abserr}{\varepsilon_{\mathrm{abs}}}
\newcommand{\relerr}{\varepsilon_{\mathrm{rel}}}
\newcommand{\maxerr}{\varepsilon_{\mathrm{max}}}
\newcommand{\meanerr}{\varepsilon_{\mathrm{mean}}}

% Pipeline notation — use inside math: $\pipeline$
\newcommand{\pipeline}{\mathrm{src}\!\to\!\mathrm{lex}\!\to\!\mathrm{parse}\!\to\!\mathrm{AST}\!\to\!\mathrm{eval}}

% Disclaimer footnote helper
\newcommand{\trademarknote}{%
  \pine{} and \tv{} are trademarks of TradingView, Inc.
  \pyne{} is an independent, unofficial implementation and is not affiliated with,
  authorized by, sponsored by, or endorsed by TradingView, Inc.
}

% Figure path helper (papers live one level under .papers/)
\graphicspath{{figures/}{../common/figures/}}


\usepackage{amsmath,amssymb,amsfonts}
\usepackage{amsthm}
\usepackage{booktabs}
\usepackage{graphicx}
\usepackage{xcolor}
\usepackage{hyperref}
\usepackage{url}
\usepackage{algorithm}
\usepackage{algpseudocode}
\usepackage{listings}
\usepackage{tikz}
\usepackage{pgfplots}
\pgfplotsset{compat=1.17}
\usetikzlibrary{arrows.meta,positioning,fit,calc,shapes.geometric,shapes.multipart,decorations.pathreplacing,backgrounds,chains,matrix,patterns}
\usepackage{subcaption}
\usepackage{multirow}
\usepackage{tabularx}
\usepackage{enumitem}
\usepackage{microtype}
\usepackage{balance}
\usepackage{csquotes}
\usepackage{natbib}

% Colors
\definecolor{pyneorange}{HTML}{F97316}
\definecolor{pyneblue}{HTML}{2563EB}
\definecolor{pynegray}{HTML}{6B7280}
\definecolor{pynegreen}{HTML}{16A34A}
\definecolor{codebg}{HTML}{F8FAFC}
\definecolor{codekw}{HTML}{7C3AED}
\definecolor{codestr}{HTML}{B45309}
\definecolor{codecomment}{HTML}{64748B}

\hypersetup{
  colorlinks=true,
  linkcolor=pyneblue,
  citecolor=pynegreen,
  urlcolor=pyneorange,
  pdftitle={},
  pdfauthor={HOOX / PYNE Project}
}

% Code listings
\lstdefinestyle{pine}{
  basicstyle=\ttfamily\footnotesize,
  backgroundcolor=\color{codebg},
  keywordstyle=\color{codekw}\bfseries,
  stringstyle=\color{codestr},
  commentstyle=\color{codecomment}\itshape,
  breaklines=true,
  frame=single,
  rulecolor=\color{black!20},
  numbers=left,
  numberstyle=\tiny\color{pynegray},
  numbersep=6pt,
  xleftmargin=1.5em,
  framexleftmargin=1.2em,
  showstringspaces=false,
  tabsize=2,
  morekeywords={indicator,strategy,library,plot,plotshape,fill,hline,var,varip,if,else,for,while,switch,true,false,na,and,or,not,import,export,type,method,enum},
  morecomment=[l]{//},
  morestring=[b]",
}
\lstdefinestyle{python}{
  language=Python,
  basicstyle=\ttfamily\footnotesize,
  backgroundcolor=\color{codebg},
  keywordstyle=\color{codekw}\bfseries,
  stringstyle=\color{codestr},
  commentstyle=\color{codecomment}\itshape,
  breaklines=true,
  frame=single,
  rulecolor=\color{black!20},
  numbers=left,
  numberstyle=\tiny\color{pynegray},
  numbersep=6pt,
  xleftmargin=1.5em,
  framexleftmargin=1.2em,
  showstringspaces=false,
  tabsize=4,
}
\lstset{style=python}

% Theorem environments
\theoremstyle{plain}
\newtheorem{theorem}{Theorem}[section]
\newtheorem{lemma}[theorem]{Lemma}
\newtheorem{proposition}[theorem]{Proposition}
\newtheorem{corollary}[theorem]{Corollary}
\theoremstyle{definition}
\newtheorem{definition}[theorem]{Definition}
\newtheorem{example}[theorem]{Example}
\newtheorem{remark}[theorem]{Remark}

% Shortcuts
\newcommand{\pyne}{\textsc{Pyne}}
\newcommand{\pine}{Pine Script\texttrademark}
\newcommand{\tv}{TradingView\textregistered}
\newcommand{\asdl}{\textsc{Asdl}}
\newcommand{\antlr}{\textsc{Antlr}4}
\newcommand{\numba}{\textsc{Numba}}
\newcommand{\lsp}{\textsc{Lsp}}
\newcommand{\ohlcv}{\textsc{Ohlcv}}
\newcommand{\udt}{\textsc{Udt}}
\newcommand{\jit}{\textsc{Jit}}
\newcommand{\api}{\textsc{Api}}
\newcommand{\cli}{\textsc{Cli}}
% AST acronym (text mode). Keep amsmath's math asterisk \ast intact.
\DeclareRobustCommand{\Ast}{\textsc{Ast}}
% Papers in this series write \ast{} for the AST acronym in prose.
% Redefine robustly; use \mathast for a binary asterisk in formulas.
\let\mathast\ast
\DeclareRobustCommand{\ast}{\textsc{Ast}}
\newcommand{\ir}{\textsc{Ir}}
\newcommand{\ta}{\texttt{ta}}
\newcommand{\code}[1]{\texttt{#1}}
\newcommand{\file}[1]{\texttt{#1}}
\newcommand{\eg}{e.g.}
\newcommand{\ie}{i.e.}
\newcommand{\etal}{\textit{et~al.}}
\newcommand{\resp}{respectively}

% Math helpers
\newcommand{\R}{\mathbb{R}}
\newcommand{\N}{\mathbb{N}}
\newcommand{\Z}{\mathbb{Z}}
% Text-mode NA token for prose; use $\namath$ inside formulas if needed.
\DeclareRobustCommand{\na}{\textsf{na}}
\newcommand{\namath}{\mathsf{na}}
\newcommand{\series}[1]{\mathbf{#1}}
\newcommand{\baridx}{i}
\newcommand{\bars}{n}
\newcommand{\period}{p}
\newcommand{\abserr}{\varepsilon_{\mathrm{abs}}}
\newcommand{\relerr}{\varepsilon_{\mathrm{rel}}}
\newcommand{\maxerr}{\varepsilon_{\mathrm{max}}}
\newcommand{\meanerr}{\varepsilon_{\mathrm{mean}}}

% Pipeline notation — use inside math: $\pipeline$
\newcommand{\pipeline}{\mathrm{src}\!\to\!\mathrm{lex}\!\to\!\mathrm{parse}\!\to\!\mathrm{AST}\!\to\!\mathrm{eval}}

% Disclaimer footnote helper
\newcommand{\trademarknote}{%
  \pine{} and \tv{} are trademarks of TradingView, Inc.
  \pyne{} is an independent, unofficial implementation and is not affiliated with,
  authorized by, sponsored by, or endorsed by TradingView, Inc.
}

% Figure path helper (papers live one level under .papers/)
\graphicspath{{figures/}{../common/figures/}}
